\section{Profundidad coinductiva de la cuatrirrelación}
\label{sec:profundidad-cuatrirrelacion-ch}

Los lectores regionales y el cierre dodecafásico publican, respectivamente,
\(\pi,\varphi,e\) y \(\alpha\) desde el estado enriquecido común. Reservamos
\(N\) para la profundidad operatoria y \(n\) para la longitud decimal:
no se identifican ambos índices ni se presupone un coste uniforme de lectura.
Para cada realización exacta \(\chi\), se conserva separadamente su parte
entera y se denota por \(p_{\chi,n}\) el prefijo de longitud \(n\) de su
parte fraccionaria, en la expansión decimal canónica no eventualmente igual
a nueve. Se define
\[
 \rho^{\rm dec}_{n+1,n}(d_1\ldots d_nd_{n+1})=d_1\ldots d_n,
 \qquad p_{\chi,0}:=\varnothing.
\]
La unicidad de \(p_{\pi,n},p_{\varphi,n},p_{e,n}\) procede de la unicidad del
carácter exacto ya generado y de su cilindro decimal canónico. La de
\(p_{\alpha,n}\) procede de la clausura entera dodecafásica, de su identidad
de colas y de la compatibilidad proyectiva probada en el
teorema~\ref{thm:alpha-salida-dodecafasica-completa}. El resolvente analítico
de orden nueve aporta una caracterización finita adicional, mientras el
lector analítico completo calcula los coeficientes restantes mediante
\(b_{10+3n}=\lambda/729^n\), prueba convergencia uniforme y construye
intervalos racionales incluidos en los cilindros enteros. De este modo
identifica esa misma sección a toda profundidad por el
teorema~\ref{prop:alpha-dos-limites}. El jet finito no se utiliza como
generador ni como sustituto de la prolongación coinductiva.

\begin{theorem}[Compatibilidad proyectiva de \(\pi,\varphi,e,\alpha\) a toda profundidad]
\label{thm:profundidad-cuatrirrelacion-ch}
La acción coinductiva del estado común determina
\[
 \begin{gathered}
  \forall\chi\in\{\pi,\varphi,e,\alpha\}\;
  \forall n\ge1\quad\exists!\,p_{\chi,n},\\
  \rho^{\rm dec}_{n+1,n}(p_{\chi,n+1})=p_{\chi,n},
  \qquad
  p_{\chi,\infty}:=(p_{\chi,n})_{n\ge0}
  \in\varprojlim_{n\ge0}\{0,\ldots,9\}^{n}.
 \end{gathered}
\tag{D1}
\label{eq:profundidad-cuatrirrelacion-ch}
\]
En particular,
\[
 (p_{\pi,\infty},p_{\varphi,\infty},p_{e,\infty},p_{\alpha,\infty})
\tag{D2}
\label{eq:cuatrirrelacion-infinita-ch}
\]
es una publicación coinductiva correlacionada del estado enriquecido único.
\end{theorem}

\begin{proof}
El resultado~\ref{gen:prefijos-regionales} fija las realizaciones exactas
de los tres lectores regionales. El
teorema~\ref{thm:alpha-salida-dodecafasica-completa} construye la sección
completa \(\alpha_{\mathrm{HMT}}\) y prueba la compatibilidad de todos sus
prefijos. La expansión decimal canónica de cada una de estas
cuatro realizaciones tiene un único prefijo de cada longitud; truncar el
prefijo de longitud \(n+1\) da el de longitud \(n\). Aplicar esos lectores
al estado común da la familia correlacionada enunciada. La existencia y
compatibilidad de esta publicación exacta no identifican \(n\) con \(N\)
ni añaden una cota de coste a los procedimientos de lectura ya establecidos.
\end{proof}

El orden constructivo interno conserva los lectores de
\(\pi,\varphi,e\), el estado firmado, el sello dodecafásico entero \(K\) y la rueda
dodecafásica que publica \(\alpha\); ninguna constante convencional se
promueve a generador. Una ejecución a mil, diez mil o un millón de cifras
verifica la instancia computacional finita correspondiente de
\eqref{eq:profundidad-cuatrirrelacion-ch}; no constituye una cota del
teorema. El teorema es la familia compatible para todo \(n\) y la sección
infinita determinada por ella.

Esta profundidad no se utiliza para inferir la hipótesis del continuo. Las cuatro publicaciones
certifican la acción coinductiva del mismo generador; el cierre cardinal ya fue
obtenido mediante la cobertura, la condensación y las dos inyecciones. Su
criterio de comprobación es doble: compatibilidad de todas las truncaciones
de \eqref{eq:profundidad-cuatrirrelacion-ch} y procedencia interna de los
datos de la rueda dodecafásica que publica \(\alpha_{\rm HMT}\).
